\section{最大熵从约束中选择最少假设的分布}
\label{sec:13-Machine-Learning/11-Information-Theoretic-Learning/02}

有人告诉你一个骰子点数 \(X\in\{1,\ldots,6\}\) 的长期平均值是 4.2。就这一条信息。让你设定每个面朝上的概率。

均匀分布？均值是 3.5，不对。那把大部分概率给 6，再随便分配剩下的？可以。按等差数列排？也可以。满足"均值=4.2"这个约束的分布在六维单纯形上有一整片。该选哪个？

最大熵原则的回答：选熵最大的那个。在所有满足约束的分布中，挑一个概率最"散"的——不把概率额外集中到任何没有数据支持的地方。它不是断言现实世界"最随机"；它是在说，已知的我已经尊重了，除此以外我不加任何多余的结构。

写成优化问题：

\[
\max_p -\sum_xp(x)\log p(x)
\]

满足：

\[
\sum_xp(x)=1,\qquad
\sum_xp(x)f_j(x)=c_j,\qquad
p(x)\ge0.
\]

\(f_j\) 是特征函数，\(c_j\) 是已知的矩约束。第一行约束让概率和为 1；第二行说"这些统计事实我确实知道"；目标函数说"除此之外什么都别假设"。

这个思想放在日常场景中很自然。电力公司出于隐私或存储限制，只保留了每小时用电量的取值范围、样本均值和超过峰值阈值的比例，原始记录删了。最大熵可以在这些汇总约束下补出一个单时段分布，粗略比较不同容量阈值下的过载概率。但补出来的分布是"在给定约束之外尽量不加结构"的分布，不是从汇总量中恢复出的真实用电规律。

这里有两根容易被忽略的线。第一，样本均值和峰值比例本身就有抽样误差。把它们当作精确等式写进模型，是把一次样本的波动固化为模型事实。实际分析可以改用允许误差的约束、对 moment mismatch 加惩罚项，并检查容量结论对约束区间是否敏感。第二，如果风险来自连续多个时段的高负载，单时段的矩根本就不包含时间依赖信息，最大熵也补不出这种依赖来——它不会无中生有。

\subsection{指数族不是凭空猜出的}

约束优化引入 Lagrange 乘子。假设最优解在可行集相对内部，对每个 \(p(x)\) 求导：

\[
-(\log p(x)+1)
+\lambda^{\trans} f(x)+\alpha=0.
\]

整理得到 \(p(x)\propto\exp\{\lambda^{\trans} f(x)\}\)，归一化后：

\[
p_\lambda(x)
=\exp\{\lambda^{\trans} f(x)-A(\lambda)\},
\]

其中 log-partition function：

\[
A(\lambda)
=\log\sum_xe^{\lambda^{\trans} f(x)}.
\]

指数族的形式不是谁拍脑袋拍出来的——它是"给定矩约束、最大化熵"这个框架的必然结果。你在上一节看到的 softmax 分类器，正是条件最大熵在这个框架下的特例。

\(A(\lambda)\) 不只是负责让概率和为 1。它对 \(\lambda\) 的导数直接给出 moments。一阶导：

\[
\frac{\partial A}{\partial\lambda_j}
=\frac{\sum_xf_j(x)e^{\lambda^{\trans} f(x)}}{\sum_xe^{\lambda^{\trans} f(x)}}
=\E_{p_\lambda}[f_j(X)].
\]

二阶导是协方差矩阵：

\[
\nabla A(\lambda)
=\E_{p_\lambda}[f(X)],
\qquad
\nabla^2A(\lambda)
=\Cov_{p_\lambda}(f(X))
\succeq0.
\]

\(A\) 是凸函数，Hessian 半正定。对偶问题：

\[
\min_\lambda A(\lambda)-\lambda^{\trans} c
\]

是凸优化。梯度恰好是模型 moment 与目标 moment 的差值——在零处停止时，约束刚好满足。

回到骰子。只约束均值时，解为 \(p_\lambda(x)\propto e^{\lambda x}\)。若目标均值 \(c=3.5\)，\(\lambda=0\) 给出均匀分布；若 \(c=4.2\)，\(\lambda>0\)，高点数获得更多概率，不再是均匀的。最大熵不是"总是均匀"——均匀只是没有任何有效矩约束时的退化结果。一旦你知道均值偏离了中点，最大熵就用指数倾斜来兼容这个信息，同时不引入更高阶的结构。

\subsection{"最少假设"仍依赖支持和基准测度}

在有限离散空间，"熵最大"的含义很清楚。到了连续情形，differential entropy \(\int -p(x)\log p(x)dx\) 会随坐标变换改变，甚至可以为负——把 \(x\) 换成 \(2x\)，entropy 就变了。所以必须明确：是相对于什么基准测度、在什么支持上讨论。

给定实数变量只有均值和方差约束时，Gaussian 最大化 differential entropy。再加一个非负支持约束，答案改变——截断指数族或截断正态可能替代 Gaussian。再加一个上界，形式又变。约束不变时换了坐标，分布也会跟着变——不是自然规律变了，是你度量的尺子变了。

约束本身也可能出问题。让骰子均值为 7.0——骰子最大才 6——可行集为空，没有分布能满足。对偶优化发散：不是求解器坏了，是约束本就矛盾。把不可行的约束当作精确等式写进模型，再好的优化器也给不出合法答案。

矩约束冗余时，\(\lambda\) 可能不唯一，但诱导的概率分布仍唯一。参数可识别性和分布可识别性是两码事。数值实现上，计算 \(A(\lambda)\) 要用稳定的 log-sum-exp——跟上一节 softmax 损失中减最大值的技巧同理。然后检查概率是否归一、moment residual 是否接近零、Hessian 是否半正定、primal-dual gap 是否闭合。这些检查不是走过场——它们验证你解出来的是否真是一个合法的最大熵分布。

\subsection{最大熵分类器怎样出现}

最大熵不限于建模 \(p(x)\)。对条件分布 \(p(y\mid x)\) 施加经验特征矩约束、最大化条件熵，得到的是 log-linear 模型：

\[
p_\lambda(y\mid x)
\propto\exp\{\lambda^{\trans} f(x,y)\}.
\]

二分类 Logistic 回归令 \(f(x,y)=y\cdot\phi(x)\)；多分类 softmax 令 \(f(x,y)\) 为 one-hot 与特征向量的外积。上一节用 cross-entropy 训练的 softmax 分类器，其函数形式正来自这个最大条件熵框架。线性 logits 不是随便选的——在"特征矩匹配、否则不加额外结构"的原则下，它是唯一合理的形式。选择它有明确的建模来源。

\subsection{约束怎样改变答案}

支持还是 \(\{1,\ldots,6\}\)。只约束均值：解属于 \(p(x)\propto e^{\lambda_1x}\)。再约束二阶矩：解扩展为 \(p(x)\propto e^{\lambda_1x+\lambda_2x^2}\)。多一个约束，可达到的最大熵只降不升，分布形状跟着变。每加一条新信息，分布就偏离"不加假设"更远一步——这正是你想要的：知道越多，分布越具体。

如果新增约束与已有约束线性相关，不增加新信息，分布不变。此时 \(\lambda\) 空间多了一维冗余，但概率唯一。关键区分是：哪些约束在改变分布形态，哪些只是在重新参数化。

最后一条警示。如果真实数据生成过程不属于指数族，最大熵分布仍可在给定矩下给出一个尽可能不加结构性假定的推断——它不是真实分布的无偏估计，只是在尊重已知矩的前提下，把未知部分填成了熵最大的形状。将这个分布外推到没有矩约束的量上——高阶矩、尾部概率、时间相关性——后果自负。那些位置的数据是真的什么也没说。
